\documentclass[10pt,a4paper]{amsart}
\usepackage[margin=19mm]{geometry}
\usepackage{amsmath,amssymb,mathtools}
\usepackage[scr=rsfs,cal=euler]{mathalfa}
\usepackage{xfrac}
\usepackage[unicode,hidelinks,bookmarks=false]{hyperref}
\newcommand{\ie}{i.e.\ }
\newcommand{\cf}{cf.\ }
\newcommand{\resp}{respectively\ }
\newcommand{\Subsection}{\subsection}
\providecommand{\nobreakdash}{\nobreak\hbox{-}}
\setlength{\emergencystretch}{2em}
\multlinegap0pt
\usepackage{amssymb}
\newtheorem{theorem}{Theorem}
\newcommand{\C}{{\mathbb{C}}}
\renewcommand{\P}{{\mathbb{P}}}
\title{A Brownian-Motion Approach to the Second Main Theorem for Meromorphic Mappings and Hypersurfaces with Truncated Counting Functions}
\author{Nguyen Linh Chi, Si Duc Quang}
\date{Source: arXiv:2605.20762v1 (CC BY 4.0)}
\begin{document}
\maketitle

\noindent\emph{Source and licence.} A Brownian-Motion Approach to the Second Main Theorem for Meromorphic Mappings and Hypersurfaces with Truncated Counting Functions. Version arXiv:2605.20762v1 (CC BY 4.0), distributed by arXiv under the Creative Commons Attribution 4.0 International licence, http://creativecommons.org/licenses/by/4.0/, as recorded in the arXiv licence metadata for this exact version and shown on the abstract page https://arxiv.org/abs/2605.20762v1. Full licence notice: \texttt{licenses/arxiv-2605.20762-CC-BY-4.0.txt}.\par\smallskip

\noindent\textit{Editorial scope.} Counted unit: the article's theorem 1.2 (source0.tex lines 159--165). The article's complete original proof of it is delivered with it (source0.tex lines 528--586, one \texttt{proof} environment of 59 lines, which the article places away from the statement). No mathematics is written, repaired or re-derived: the statement and the proof are the article's own lines, in the article's own notation, and no retained line is substituted or re-flowed.  Retained uncounted as the source-local context the proof needs: lines 150--153.  Nothing was omitted from any retained span, so the package carries no omission record.  No retained line contains a citation, so the wrapper carries no bibliography and no bibliography entry.  No retained line contains a figure environment, an image-inclusion command or drawing code, so no image is delivered and the package declares no asset dependency.  Editorial formatting only, in addition: an AMS \texttt{amsart} preamble that restates the article's own declarations and macros used by the retained lines, namely the environments \texttt{theorem} on separate counters, together with the standard packages those lines need.  The retained environments print in this document as Theorem 1 in order of appearance, whereas the article prints the same items with the numbering of its own full document.  The complete native source of the article as distributed in the arXiv e-print for this exact version is bundled unchanged as \texttt{sources/201015/source0.tex}.
\begin{theorem}\label{1.1}
Let $V$ be a projective subvariety of $\P^n(\C)$ and let $\{D_1,\ldots,D_q\}$ be a family of $q$ hypersurfaces of $\P^n(\C)$ which has distributive constant $\Delta$ with respect to $V$. Let $M=H_V(d)-1$ and $d=lcm(\deg D_1,\ldots,\deg D_q)$. Let $f$ be a holomorphic map from $\C$ into $V$ which is nondegenerate over $I_d(V)$. Then, for any positive $\epsilon$ and $\delta$, we have
$$\|\,(q-\Delta(M+1+\epsilon))T_f(r)\le\sum_{j=1}^q\frac{1}{d}N^{[M]}_f(r,D_j)+\Delta\delta \log r + O(1).$$
\end{theorem}

\begin{theorem}\label{1.2}
Let $V$ be a projective subvariety of $\mathbb P^n(\mathbb C)$.  Let $\{D_i\}_{i=1}^q$ be a family of $q$ hypersurfaces in $\mathbb P^n(\mathbb C)$ with the distributive constant $\Delta$ with respect to $V$. Let $d=lcm (deg D_1,...,deg D_q)$. Let $f$ and $g$ be holomorphic maps from $\C$ into $V$ which are nondegenerate over $I_d(V)$.  Assume that $f=g \text{ on }\bigcup_{i=1}^q(f^{-1}(D_i)\cup g^{-1}(D_i)).$
\begin{itemize}
    \item[(a)] If $q>\Delta\left(\frac{2k(H_V(d)-1)}{d}+ H_V(d)\right),$ then $f=g.$
    \item[(b)] Suppose further that $f^{-1}(D_i\cap D_j)=\varnothing$ for every $1\le i< j\le q.$ If $q>\frac{2(H_V(d)-1)}{d}+\Delta H_V(d),$ then $f=g.$
\end{itemize}
\end{theorem}

\begin{proof}[Proof of Theorem \ref{1.2}]
Assume that $\tilde f=(f_0,\ldots,f_n)$ and $\tilde g=(g_0,\ldots,g_n)$ are reduced representations of $f$ and $g,$ respectively. Let $Q_1,\ldots,Q_q$ be the homogeneous polynomials defining$D_1,\ldots,D_q$ respectively. Without loss of generality, we may assume that $\deg Q_j=d$ for all $1\le j\le q$.

Suppose that $f\ne g$. Then there exist two indices $s,t$ with $0\le s<t\le n$ such that $ H:=f_sg_t-f_tg_s\not\equiv 0.$
By the definition of the characteristic function and by the Jensen formula, we have
\begin{align}\label{5}
\begin{split}
N_H(r)&=\int_{0}^{2\pi}\log |f_s(re^{i\theta})g_t(re^{i\theta})-f_t(re^{i\theta})g_s(re^{i\theta})|\frac{d\theta}{2\pi}\\ 
& \le \int_{0}^{2\pi}\log \|f(re^{i\theta})\|\frac{d\theta}{2\pi} + \int_{0}^{2\pi}\log \|g(re^{i\theta})\|\frac{d\theta}{2\pi} \\
&=T_f(r)+T_g(r).
\end{split}
\end{align}

(a) By the assumption of the theorem, we have $H=0$ on $\bigcup_{i=1}^qf^{-1}(D_i)\cup g^{-1}(D_i)$. 
If $z$ is a zero of exactly $\ell$ functions $Q(\tilde f)$ then $\Delta\ge\frac{\ell}{k}$, i.e., $\ell\le\Delta k.$
Therefore, 
$$\nu^0_H(z)\ge \Delta k\sum_{i=1}^q\min\{1,\nu^0_{Q_i(f)}\}.$$
It follows that
\begin{align}\label{6}
N_H(r)\ge \Delta k\sum_{i=1}^qN^{[1]}_{Q_i(f)}(r).
\end{align}
Combining (\ref{5}) and (\ref{6}), we have
$$ T_f(r)+T_g(r) \ge \Delta k\sum_{i=1}^qN^{[1]}_{Q_i(f)}(r).$$
Similarly, we have
$$ T_f(r)+T_g(r) \ge \Delta k\sum_{i=1}^qN^{[1]}_{Q_i(g)}(r).$$
Summing both sides of the above inequalities, we get
\begin{align*}
 2(T_f(r)+T_g(r))& \ge \Delta k\left(\sum_{i=1}^qN^{[1]}_{Q_i(f)}(r)+\sum_{i=1}^qN^{[1]}_{Q_i(g)}(r)\right).
\end{align*}
This yields that
\begin{align*}
\frac{2}{\Delta k}(T_f(r)+T_g(r))&\ge\sum_{i=1}^qN^{[1]}_{Q_i(f)}(r)+\sum_{i=1}^qN^{[1]}_{Q_i(g)}(r)\\
&\ge \sum_{i=1}^q\dfrac{1}{H_V(d)-1}N^{[H_V(d)-1]}_{Q_i(f)}(r)+\sum_{i=1}^q\dfrac{1}{H_V(d)-1}N^{[H_V(d)-1]}_{Q_i(g)}(r)\\
&\ge \dfrac{d}{H_V(d)-1}(q-\Delta (H_V(d)+\epsilon))(T_f(r)+T_g(r))+o(T_f(r)+T_g(r)).
\end{align*}
Letting $r\longrightarrow +\infty$ and then letting $\epsilon\longrightarrow 0$, we get 
$$2\Delta k\ge \dfrac{d}{H_V(d)-1}(q-\Delta H_V(d)),$$
$$\text{i.e., } q\le \Delta\left(\frac{2k(H_V(d)-1)}{d}+ H_V(d)\right).$$
This is a contradiction.
Hence $f=g$.

(b) Since $H=0$ on $\bigcup_{i=1}^qf^{-1}(D_i)\cup g^{-1}(D_i)$, we have
$$\nu^0_H(z)\ge \sum_{i=1}^q\min\{1,\nu^0_{Q_i(f)}\}.$$
It follows that
$$N_H(r)\ge \sum_{i=1}^qN^{[1]}_{Q_i(f)}(r).$$
Combining (\ref{5}) and this inequality, we obtain
$$ T_f(r)+T_g(r) \ge \sum_{i=1}^qN^{[1]}_{Q_i(f)}(r)\, \text{ and }\, T_f(r)+T_g(r) \ge \sum_{i=1}^qN^{[1]}_{Q_i(g)}(r).$$
Summing-up both sides of the above two inequalities and applying Theorem \ref{1.1} we have
\begin{align*}
2(T_f(r)+T_g(r))&\ge\sum_{i=1}^qN^{[1]}_{Q_i(f)}(r)+\sum_{i=1}^qN^{[1]}_{Q_i(g)}(r)\\
&\ge \sum_{i=1}^q\dfrac{1}{H_V(d)-1}N^{[H_V(d)-1]}_{Q_i(f)}(r)+\sum_{i=1}^q\dfrac{1}{H_V(d)-1}N^{[H_V(d)-1]}_{Q_i(g)}(r)\\
&\ge \dfrac{d}{H_V(d)-1}(q-\Delta (H_V(d)+\epsilon))(T_f(r)+T_g(r))+o(T_f(r)+T_g(r)).
\end{align*}
Letting $r\longrightarrow +\infty$ and then letting $\epsilon\longrightarrow 0$, we get 
$$2 \ge \dfrac{d}{H_V(d)-1}(q-\Delta H_V(d)),$$
$$\text{i.e., } q\le \frac{2(H_V(d)-1)}{d}+\Delta H_V(d).$$
This is a contradiction.
Hence $f=g$.
\end{proof}

\end{document}
